\chapter{Command reference}
\label{app:cheatsheet}
\begingroup
\lstset{basicstyle=\ttfamily\footnotesize}

Commands assume the environment of the package named in the heading
(\cmd{pixi cvmfs} in its checkout), \file{work/} as the current directory
except Geometry build commands, which run in its checkout, and the geometry
tutorial commands, which run in \file{\$SHIP\_TUTORIAL/geometry-tutorial};
\texttt{\$CODE} points to
\file{framework-booklet/code} in the supplied archive.

\section*{git \hfill\normalfont\small(\cref{ch:git,ch:branches,ch:remotes}; full table in \cref{sec:gitcheatsheet})}
\begin{shell}
git switch main && git pull upstream main       # start from the latest main
git switch -c task-NN-description               # one branch per change
git add FILE && git commit -m "feat(scope): summary"
git push -u origin task-NN-description          # then open the pull request
git fetch upstream && git rebase upstream/main  # PR out of date
git push --force-with-lease                     # after a rebase (your branch only)
\end{shell}

\section*{pixi \hfill\normalfont\small(\cref{ch:pixi})}
\begin{shell}[basicstyle=\ttfamily\scriptsize]
# once, in ~/.bashrc (exercise 4.1):
export PATH=/cvmfs/ship-nightlies.cern.ch/pixi/0.72.0-rattlerfs-layers-d5a7568/bin:$PATH
\end{shell}
\begin{shell}
# a new workspace:
pixi init && pixi add --no-install python ruff
pixi cvmfs <task or command>      # one command in the environment
pixi cvmfs                        # a shell in the environment; exit to leave
pixi task add <name> "<command>"  # name a command
pixi lock                         # after editing pixi.toml: commit both
pixi update --no-install          # newer versions: commit the new pixi.lock
pixi clean                        # remove a .pixi/ installed by mistake
fs lq ~                           # AFS quota and usage
\end{shell}
On lxplus, never \cmd{pixi run}, \cmd{pixi shell} or \cmd{pixi install}:
they install into \file{.pixi/} on AFS.

\section*{CI \hfill\normalfont\small(\cref{ch:ci})}
\begin{shell}
pixi cvmfs <task>  # what CI runs as: pixi run --locked <task>
pixi cvmfs lint    # what "Lint passed" checks, before you push
\end{shell}

\section*{Team analysis \hfill\normalfont\small(\cref{ch:team}, in the team repository)}
\begin{shell}
pixi cvmfs check-task NN  # test your task alone
pixi cvmfs check          # test every task (as CI does)
pixi cvmfs analysis       # run the chain as far as it goes
\end{shell}

\section*{Geometry \hfill\normalfont\small(\repo{Geometry})}
\begin{shell}
pixi cvmfs geometry                                   # -> build/ship_geometry.db
pixi cvmfs ./build/apps/validate_geometry build/ship_geometry.db
pixi cvmfs gdml                                       # -> build/ship_geometry.gdml
pixi cvmfs ./build/apps/build_geometry --list         # the subsystem names
pixi cvmfs ./build/apps/build_geometry --standalone NAME -o NAME.db
\end{shell}

\section*{Geometry tutorial and viewer \hfill\normalfont\small(\file{\$CODE/geometry})}
\begin{shell}
pixi cvmfs plane                       # -> calo_plane.db (GeoModel from scratch)
pixi cvmfs view [file.db | macro.mac]  # Geant4 window; default calo_plane.mac
pixi cvmfs check                       # tree, mass, overlaps; no window
pixi cvmfs picture                     # -> geometry.png; no window
pixi cvmfs view ship.mac               # SHiP, one subsystem at a time
\end{shell}
In the viewer: \cmd{/ship/geometry/open F.db} (before \cmd{/run/initialize}),
\cmd{/vis/drawVolume NAME}, \cmd{/vis/scene/add/volume NAME},
\cmd{/vis/drawTree NAME}, \cmd{/geometry/test/run}, \cmd{/gun/particle P},
\cmd{/run/beamOn N}.

\section*{Field maps \hfill\normalfont\small(\repo{field\_service})}
\begin{shell}
fieldmap_dump MAP.root                      # metadata
echo "0 0 0" | fieldmap_dump MAP.root NAME  # evaluate at points (mm -> T)
plot_field_map MAP.root NAME plots.root     # B_y plots, ROOT + PDF
generate_constant_fieldmap OUT.root NAME Bx By Bz xmin xmax ymin ymax zmin zmax
fairship_to_fieldmap [--symmetry quadrant_dipole] IN.root OUT.root
\end{shell}

\section*{Simulation \hfill\normalfont\small(\repo{aegir})}
\begin{shell}
POT=1000                                 # protons on target, stage 1
phlex -j 4 -c <(jsonnet --tla-code events=$POT --tla-code seed=42 \
    [--tla-code ecut=0.5] [--tla-str output=ft_production] \
    $CODE/workflows/ft_production.jsonnet)
python $CODE/analysis/plane_summary.py ft_production.root --pot $POT
phlex -j 4 -c <(jsonnet [--tla-code events=100] [--tla-code pdg=-13] \
    [--tla-code seed=42] [--tla-code field=false] [--tla-str output=muon_gun] \
    $CODE/workflows/muon_gun.jsonnet)       # artificial muons from z = 27 m
python $CODE/analysis/sim_summary.py muon_gun.root [other.root]
N=$(python "$SHIP_TUTORIAL/aegir/scripts/count_entries.py" \
      ft_production.root)                  # optional: replay stage 1
phlex -j 4 -c <(jsonnet --tla-str infile=ft_production.root --tla-code events=$N \
    $CODE/workflows/ft_replay.jsonnet)
\end{shell}

\section*{Event display \hfill\normalfont\small(\repo{sea-cucumber}, run in its checkout)}
\begin{shell}
pixi cvmfs run --geometry DB --data muon_gun.root --view views/default.toml --event I
pixi cvmfs run --geometry DB --inspect 4 [--inspect-match "*trackers*"]
pixi cvmfs geo-cache --geometry DB --view views/default.toml
pixi cvmfs web-data --geometry DB --data muon_gun.root --view views/default.toml \
                    --events all
PORT=$((8000 + $(id -u) % 1000)); hostname; echo $PORT
pixi cvmfs python -m http.server $PORT --bind 127.0.0.1 --directory web
# on your own machine, then open http://localhost:8080:
ssh -N -L 8080:localhost:<PORT> <you>@lxplusNNN.cern.ch
\end{shell}

\section*{Digitisation \hfill\normalfont\small(\repo{Shannon})}
\begin{shell}
phlex -j 4 -c <(jsonnet --tla-str infile=muon_gun.root --tla-code entries=100 \
    --tla-code pot=1000 [--tla-str output=ft_digi] $CODE/workflows/ft_digi.jsonnet)
python $CODE/analysis/digi_residuals.py ft_digi.root --collection straw_tubes_hits
\end{shell}

\section*{Analysis}
\begin{shell}
cp $CODE/analysis/first_sim.py $CODE/analysis/first_digi.py .   # to extend
python first_sim.py muon_gun.root          # hits per detector (aegir file)
python first_digi.py ft_digi.root          # hits per collection (Shannon file)
python $CODE/analysis/muon_background.py muon_gun.root [--output PREFIX]
python $CODE/analysis/muon_background.py ft_replay.root --pot $POT   # a replay
python $CODE/analysis/solutions/spectrometer_charge.py muon_gun.root muon_gun_mup.root
\end{shell}

\section*{Looking inside any file}
\begin{py}
import ROOT
r = ROOT.RNTupleReader.Open("events", "muon_gun.root")    # or "straw_tubes_hits", ...
r.PrintInfo(); r.Show(0)
\end{py}
\endgroup
